% Propietario conservado: /Users/ruben/Documents/New project/output/REVISION_ARGUMENTAL_APERTURAS_CIERRES_20260915/VII/delivery/MANUSCRIPT_VII_ES_EN_COMPLETE/source_es/manuscrito/30e_corriente_extension_minima.tex
% Artículo VII. Incorporación aditiva, 10-09-2026.
% RESULTADO_RECUPERADO: fuentes_conservadas/18_clausura_energetica.tex,
% líneas 357--485: Gramiano, peso, extensión mínima y eliminación de Schur;
% líneas 149--180 y 235--295: normalización y unicidad de la respuesta.
% Residencia anterior en VII: 30_pantalla_y_respuesta.tex,
% vii:prop:incidencia, vii:eq:lambda-pantalla, vii:prop:minima-extension.
% FORMALIZACION_NUEVA / CERTIFICADO_NUEVO: diferencial del mínimo con
% incidencia, frontera y normalización variables; incidencia transportada
% por bloques y composición con coordenadas de conexión.
% La familia de rutas que realiza cada incidencia se mantiene explícita.
% No se atribuyen al propietario histórico ni esa familia diferenciable
% completa ni la evaluación de una amplitud material de espín.
% Dependencias: 30, 30c y la convención variacional de 31. No modifica
% esos fragmentos ni constituye por sí solo una realización material.

\section{Différentielle de l'extension minimale de frontière}
\label{vii:vii:sec:corriente-extension-minima}

L'incidence cubique détermine simultanément le champ d'extension
minimale et le poids de sa réponse. Pour une donnée de frontière, tous deux sont
calculés au moyen de l'inverse du gramien d'incidence. On développe
maintenant la différentielle de cette même opération, en conservant les variations
de la frontière, du transport et de l'échelle de normalisation. De cette
manière, le champ qui intervient dans la réponse est sélectionné par
le problème variationnel d'écran.

\subsection{Champ minimisant et poids de réponse}
\label{vii:vii:subsec:extension-minima-parametrizada}

Soient \(H\) et \(Q\) des espaces hermitiens de dimension finie, à produits fixes et antilinéaires en la première variable. Le cas réel s'obtient en remplaçant les adjoints par les transposées. Dans un ouvert de paramètres réels \(\Omega\), considérons des familles de classe \(C^1\)
\[
 \mathsf C(\theta):H\longrightarrow Q,\qquad
 b(\theta)\in Q,\qquad \chi(\theta)>0,
 \qquad \operatorname{Ran}\mathsf C(\theta)=Q.
\]
La surjectivité équivaut à un rang en lignes maximal dans des bases orthonormées. Nous définissons le gramien, le potentiel conjugué et l'extension
\begin{equation}
 \mathsf L=\mathsf C\mathsf C^*,\qquad
 y=\mathsf L^{-1}b,\qquad f_{\min}=\mathsf C^*y,
 \qquad \Lambda=\chi\mathsf L^{-1}.
 \label{vii:vii:eq:gramiano-extension-minima}
\end{equation}
L'incidence \(\mathsf C\) et l'accouplement cellulaire \(M\) de \eqref{vii:vii:eq:coordenadas-corriente-material} sont des opérateurs différents.

\begin{proposition}[Sélection variationnelle de l'extension]
\label{vii:vii:prop:minimo-parametrizado}
L'extension \(f_{\min}\) est l'unique minimisante de la norme parmi
les champs \(f\in H\) qui satisfont \(\mathsf C f=b\). Son énergie est
\begin{equation}
 E(\theta)=\chi\min_{\mathsf C f=b}\|f\|^2
 =\chi\langle b,\mathsf L^{-1}b\rangle
 =\langle b,\Lambda b\rangle.
 \label{vii:vii:eq:energia-extension-minima}
\end{equation}
Le champ \(f_{\min}\), le potentiel \(y\) et l'énergie \(E\)
dépendent des paramètres de manière \(C^1\).
\end{proposition}

\begin{proof}
La surjectivité de \(\mathsf C\) rend \(\mathsf C^*\) injectif, de sorte que \(\langle z,\mathsf Lz\rangle=\|\mathsf C^*z\|^2>0\) pour \(z\ne0\). Ainsi, \(\mathsf L\) est inversible et \(\mathsf C f_{\min}=\mathsf L y=b\). Toute solution est \(f_{\min}+z\), avec \(z\in\ker\mathsf C\). Puisque \(f_{\min}\in\operatorname{Ran}\mathsf C^*=(\ker\mathsf C)^\perp\),
\[
 \|f_{\min}+z\|^2=\|f_{\min}\|^2+\|z\|^2,
 \qquad
 \|f_{\min}\|^2=\langle y,\mathsf L y\rangle
 =\langle b,\mathsf L^{-1}b\rangle.
\]
Ces identités prouvent le minimum, son unicité et sa valeur. L'inversion
est une opération différentiable sur l'ouvert des matrices inversibles ;
les expressions de \eqref{vii:vii:eq:gramiano-extension-minima} fournissent
la régularité affirmée.
\end{proof}

La réalisation cubique de \ref{vii:vii:prop:incidencia} correspond à
\(H=\mathbb R^{\{-1,1\}^3}\), au quotient d'écran
\(Q=P_\square\mathbb R^6\) et à l'incidence \(\mathsf C=M:H\to Q\).
Son gramien et sa normalisation sont
\begin{equation}
 \mathsf L_\square=12P_u+4P_-,\qquad
 \chi_m=\frac{112}{81\ell_m^2},\qquad
 \Lambda_m=\frac{28}{243\ell_m^2}(P_u+3P_-).
 \label{vii:vii:eq:normalizacion-minimo-cubico}
\end{equation}
Le facteur \(112/81\) provient de la soudure et du défaut quadratique de compression de \eqref{vii:vii:eq:lambda-pantalla}. Si \(b\) a la dimension d'une longueur, \(y\) et \(f_{\min}\) l'ont également, l'incidence est sans dimension et \(E\) est sans dimension. La conversion en action utilise la normalisation dans la droite d'action spécifiée plus loin.

\subsection{Différentielle complète du minimum}
\label{vii:vii:subsec:diferencial-minimo}

\begin{theorem}[Réponse à la variation d'incidence et de frontière]
\label{vii:vii:thm:diferencial-minimo}
Pour toute direction tangente réelle \(\delta\theta\), on a
\begin{equation}
 \delta E
 =\delta\chi\,\langle b,y\rangle
   +2\chi\operatorname{Re}
       \langle y,\delta b-(\delta\mathsf C)f_{\min}\rangle.
 \label{vii:vii:eq:diferencial-minimo-completo}
\end{equation}
Si \(\chi=112/(81\ell_m^2)\), avec \(\ell_m>0\) variable,
\begin{equation}
 \delta E=-2E\frac{\delta\ell_m}{\ell_m}
 +2\chi_m\operatorname{Re}
       \langle y,\delta b-(\delta\mathsf C)f_{\min}\rangle.
 \label{vii:vii:eq:diferencial-minimo-escala}
\end{equation}
\end{theorem}

\begin{proof}
Différentier \(\mathsf L\mathsf L^{-1}=I_Q\) donne
\[
 \delta\mathsf L^{-1}
 =-\mathsf L^{-1}(\delta\mathsf L)\mathsf L^{-1},\qquad
 \delta\mathsf L=(\delta\mathsf C)\mathsf C^*
                      +\mathsf C(\delta\mathsf C)^*.
\]
En appliquant la règle du produit à \eqref{vii:vii:eq:energia-extension-minima} et en utilisant le caractère autoadjoint de \(\mathsf L^{-1}\), il vient
\[
 \delta E=\delta\chi\,\langle b,y\rangle
 +2\chi\operatorname{Re}\langle y,\delta b\rangle
 -\chi\langle y,(\delta\mathsf L)y\rangle.
\]
Les deux termes du dernier produit sont conjugués entre eux et
\(\mathsf C^*y=f_{\min}\) ; par conséquent,
\[
 \langle y,(\delta\mathsf L)y\rangle
 =2\operatorname{Re}\langle y,(\delta\mathsf C)f_{\min}\rangle.
\]
On obtient \eqref{vii:vii:eq:diferencial-minimo-completo}. Enfin, \(\delta\chi_m=-2\chi_m\delta\ell_m/\ell_m\) produit \eqref{vii:vii:eq:diferencial-minimo-escala}.
\end{proof}

La formule incorpore la variation du champ minimisant au moyen de l'
équation qui le détermine. Elle se vérifie aussi directement :
\(\mathsf C\delta f_{\min}=\delta b-(\delta\mathsf C)f_{\min}\)
et \(f_{\min}=\mathsf C^*y\) impliquent
\[
 \delta\|f_{\min}\|^2
 =2\operatorname{Re}\langle f_{\min},\delta f_{\min}\rangle
 =2\operatorname{Re}
     \langle y,\delta b-(\delta\mathsf C)f_{\min}\rangle.
\]
Par exemple, si \(\mathsf C\) est fixe et \(b=\ell_m b_0\), avec
\(b_0\) fixe, les deux termes de
\eqref{vii:vii:eq:diferencial-minimo-escala} se compensent et l'énergie
reste constante. Ce contrôle conserve conjointement la dimension
de frontière et sa normalisation.

\subsection{Quotient fixe et changements de repère}
\label{vii:vii:subsec:marcos-minimo}

Les dérivées précédentes sont prises dans des espaces \(H,Q\) et avec des produits hermitiens fixes. Dans la présentation à six faces, l'inverse agit toujours sur \(Q=P_\square H_F\). Pour une famille de quotients \(Q(\theta)\), on fixe d'abord une trivialisation isométrique \(V_Q(\theta):Q\to Q(\theta)\), et de même pour le domaine. Les dérivées de ces applications sont incluses lors de la différentiation des expressions transportées vers les espaces fixes. Les formules s'appliquent dans chaque ouvert où le rang est conservé. Les changements du produit énergétique exigent de différentier également ce produit ; une métrique lorentzienne et le produit positif utilisé pour minimiser conservent ici leurs types distincts.

\begin{proposition}[Covariance du minimum]
\label{vii:vii:prop:covariancia-minimo}
Soient \(g_H,g_Q\) unitaires. Sous
\[
 \mathsf C'=g_Q\mathsf C g_H^*,\qquad b'=g_Qb,\qquad \chi'=\chi,
\]
on a \(y'=g_Qy\), \(f'_{\min}=g_Hf_{\min}\) et \(E'=E\).
Lorsque les changements de cadre dépendent du paramètre, la différentielle
complète conserve cette égalité. En particulier, une variation pure de
cadre a \(\delta E=0\).
\end{proposition}

\begin{proof}
Le gramien transformé est \(\mathsf L'=g_Q\mathsf Lg_Q^*\),
d'où l'on obtient les transformations de \(y,f_{\min}\) et
l'égalité des énergies. Pour une variation pure de repère, en un point
avec des repères identité, soient \(A_Q^*=-A_Q\) et \(A_H^*=-A_H\)
leurs générateurs. Alors
\[
 \delta\mathsf C=A_Q\mathsf C-\mathsf C A_H,\qquad
 \delta b=A_Qb,\qquad
 \delta b-(\delta\mathsf C)f_{\min}=\mathsf C A_Hf_{\min}.
\]
Sa contribution à \eqref{vii:vii:eq:diferencial-minimo-completo} est
\(2\chi\operatorname{Re}\langle f_{\min},A_Hf_{\min}\rangle=0\).
La preuve inclut ainsi les termes dus aux cadres mobiles.
\end{proof}

\subsection{Incidence transportée et composantes de connexion}
\label{vii:vii:subsec:incidencia-conexion-minima}

La dépendance de connexion est spécifiée avant d'effectuer la
contraction. Une réalisation explicite de l'incidence par blocs
s'obtient comme suit. Soient \(\mathcal K\) une fibre hermitienne commune,
\(F=\{(i,\epsilon):1\leq i\leq3,\ \epsilon=\pm1\}\) et
\(V=\{-1,1\}^3\). Pour chaque incidence \(\nu_i=\epsilon\),
la réalisation utilisée déclare une route \(\gamma_{i,\epsilon,\nu}\)
de la fibre du sommet à celle de la face, avec le transport
\(U_{i,\epsilon,\nu}\). Les trivialisations identifient ces fibres
à \(\mathcal K\). Nous définissons
\begin{equation}
 \begin{split}
 (\mathsf B_Uf)_{i,\epsilon}
   &=\sum_{\nu_i=\epsilon}U_{i,\epsilon,\nu}f_\nu,\\
 \mathsf C_U&=(q_\square\otimes I_{\mathcal K})\mathsf B_U:
   \bigoplus_{\nu\in V}\mathcal K\longrightarrow Q\otimes\mathcal K,
 \qquad q_\square=P_\square:H_F\to Q.
 \end{split}
 \label{vii:vii:eq:incidencia-transportada-minima}
\end{equation}
Il s'agit d'une continuation variationnelle explicite de l'incidence cubique
dans la représentation des routes déclarée. Pour les transports identité,
\(\mathsf C_U=M\otimes I_{\mathcal K}\) et l'on retrouve
\eqref{vii:vii:eq:normalizacion-minimo-cubico}, tensorisée avec la fibre.
Le gramien reste positif dans un voisinage de cette réalisation,
car la positivité définie est une condition ouverte. Dans chaque domaine
de rang plein, le poids et le champ sont fixés par
\eqref{vii:vii:eq:gramiano-extension-minima}. Les routes qui réalisent les
incidences, leurs orientations et leurs mémoires sont des données de cette
réalisation du transport APP--TRIT--TPK et accompagnent la construction.

Si \(\omega=(\omega_j)\) sont des coordonnées réelles de connexion, la dérivée de \(\mathsf C_U\) se calcule en substituant dans chaque bloc
\begin{equation}
 \partial_j U_\gamma
 =\sum_{k=1}^{N}
    U_N\cdots U_{k+1}(\partial_jU_k)U_{k-1}\cdots U_1,
 \qquad U_\gamma=U_N\cdots U_1.
 \label{vii:vii:eq:derivada-incidencia-rutas}
\end{equation}
Les produits vides sont des identités. La formule provient de la règle du produit et conserve la place de chaque opération dans la route. Ainsi, \(\mathsf C_j:=\partial_j\mathsf C_U\) s'obtient à partir d'une somme finie de matrices connues dans la réalisation choisie, suivie de la projection \(q_\square\otimes I\). Si un quotient mobile est utilisé, sa dérivée est incorporée conformément à \ref{vii:vii:subsec:marcos-minimo}. La matrice cubique entière fixe correspond au cas \(\mathsf C_j=0\) ; la réponse à la connexion provient de la dépendance des transports effectivement déclarée dans \eqref{vii:vii:eq:incidencia-transportada-minima}.

\begin{theorem}[Composantes variationnelles de l'extension minimale]
\label{vii:vii:thm:componentes-minimo-conexion}
Pour l'action d'écran
\(S_{\min}=\mathfrak a_m E\), où \(\mathfrak a_m\) est sa
normalisation sur la droite d'action, écrivons
\[
 b_j=\partial_jb,\qquad \chi_j=\partial_j\chi,
 \qquad a_j=\partial_j\mathfrak a_m.
\]
Avec la convention \(\delta S_{\min}=-\tfrac12\sum_j s_j\delta\omega_j\),
les composantes sont
\begin{equation}
 \begin{split}
 s_j={}&4\mathfrak a_m\chi\operatorname{Re}
             \langle y,\mathsf C_jf_{\min}-b_j\rangle\\
      &-2\mathfrak a_m\chi_j\langle b,y\rangle-2a_jE.
 \end{split}
 \label{vii:vii:eq:componentes-minimo-conexion}
\end{equation}
Pour la normalisation cubique, le second terme équivaut à
\(4\mathfrak a_m E\,\partial_j\ell_m/\ell_m\).
Si l'appariement cellulaire de cette même action a une matrice
inversible \(\mathsf M_{\rm cel}\), ses coordonnées de courant
sont \(\sigma=\mathsf M_{\rm cel}^{-1}s\).
\end{theorem}

\begin{proof}
La règle du produit donne
\(\partial_jS_{\min}=a_jE+\mathfrak a_m\partial_jE\).
Substituer \eqref{vii:vii:eq:diferencial-minimo-completo} et multiplier
par \(-2\) prouve \eqref{vii:vii:eq:componentes-minimo-conexion}.
La dérivée de \(\chi_m\) produit l'expression d'échelle.
Enfin, comparer
\(\delta S_{\min}=-\tfrac12\delta\omega^{\mathsf T}s\)
avec \(-\tfrac12\delta\omega^{\mathsf T}\mathsf M_{\rm cel}\sigma\)
pour toute variation donne \(\mathsf M_{\rm cel}\sigma=s\),
et son inverse détermine les coordonnées.
\end{proof}

La construction retrouve un champ et un poids sélectionnés conjointement
par l'incidence, et évalue leur réponse aux variations de la même
réalisation. La donnée de frontière \(b\), la famille de routes, l'échelle
\(\ell_m\), la normalisation \(\mathfrak a_m\) et l'appariement
cellulaire conservent leurs règles de publication. L'interprétation comme
courant matériel utilise ces objets et l'action complète. La
réponse de Cartan et sa densité métrique ultérieure se composent avec
elle dans le domaine de matière correspondant, en conservant toute
dépendance du courant vis-à-vis de la connexion.
